\chapter*{Resumen}
\addcontentsline{toc}{chapter}{Resumen}

Los captadores pasivos de agua atmosf\'erica de zonas costeras semi\'aridas como Cabo de la Vela, en La Guajira, Colombia, sufren incrustaci\'on salina, que los sistemas convencionales limpian con pulsos t\'ermicos programados por un temporizador fijo. Este trabajo eval\'ua por simulaci\'on un motor de seis reglas deterministas, primera etapa simb\'olica de una arquitectura neuro-simb\'olica, frente a un temporizador de dos horas en una red de ocho nodos durante un ciclo diario completo. El estudio se implementa en el artefacto de software abierto \emph{Axioma-H Engine}, verificado mediante pruebas automatizadas e integraci\'on continua.

La respuesta es condicional. En el modelo, el control por reglas reduce la degradaci\'on acumulada de la membrana entre 61.5\,\% y 100\,\% frente al temporizador fijo al variar el flujo de niebla, la eficacia del pulso, la fatiga y la degradaci\'on t\'ermica. Frente a un temporizador con el mismo corte t\'ermico y la misma pausa ecol\'ogica, la reducci\'on en el barrido de salinidad es de 40.6\,\% a 90.9\,\%, con una excepci\'on en la sequ\'ia cuando el intervalo es de dos horas. La ventaja en recuperaci\'on depende de la incrustaci\'on diaria: aparece sobre una banda de transici\'on que se conserva con semillas fuera de muestra y se desplaza hacia arriba con la conductancia media en el tiempo; con flujos calibrados ocurre en 3 de 62 muestras de incertidumbre. Las salvaguardas eliminan el sobrecalentamiento y las violaciones ecol\'ogicas. Con flujos calibrados, abastecer a cincuenta personas al nivel b\'asico de la OMS exige entre 100 y 332 nodos, e igualar la entrega mensual por carrotanque del ICBF, entre 3 y 13.

\medskip
\noindent\textbf{Palabras clave:} captaci\'on de agua atmosf\'erica, control basado en reglas, arquitectura neuro-simb\'olica, incrustaci\'on salina, regeneraci\'on t\'ermica, La Guajira.

\chapter*{Abstract}
\addcontentsline{toc}{chapter}{Abstract}

Passive atmospheric water harvesters in coastal semi-arid regions such as Cabo de la Vela, in La Guajira, Colombia, suffer salt fouling, which conventional systems clean with thermal pulses scheduled by a fixed timer, without measuring humidity, impedance or energy reserves. This work evaluates by simulation an engine of six deterministic rules, the first symbolic stage of a neuro-symbolic architecture, against a two-hour timer in an eight-node network over a full daily cycle. The study is implemented in the open-source software artifact \emph{Axioma-H Engine}, verified through automated tests and continuous integration.

The answer is conditional. In the model, rule-based control reduces cumulative membrane degradation by 61.5\,\% to 100\,\% compared with the fixed timer across variations in fog flux, pulse efficacy, pulse fatigue and thermal degradation. Against a timer with the same thermal cut-off and ecological pause, the reduction in the salinity sweep ranges from 40.6\,\% to 90.9\,\%, with one exception under drought for a two-hour interval. The recovery advantage depends on daily salt incrustation: it appears above a transition band that holds under out-of-sample seeds and shifts upward when measured by time-averaged conductance; under calibrated fog fluxes it occurs in 3 of 62 uncertainty samples. The safeguards eliminate overheating and ecological violations. Under calibrated fluxes, supplying fifty people at the WHO basic access level requires 100 to 332 nodes, and matching the monthly water-truck delivery reported by ICBF requires 3 to 13.

\medskip
\noindent\textbf{Keywords:} atmospheric water harvesting, rule-based control, neuro-symbolic architecture, salt fouling, thermal regeneration, La Guajira.

\chapter*{Glosario de siglas y abreviaturas}
\addcontentsline{toc}{chapter}{Glosario de siglas y abreviaturas}

\begin{tabularx}{\textwidth}{@{}L{2.5cm}Y@{}}
\toprule
\textbf{Sigla} & \textbf{Significado y uso en el informe} \\
\midrule
API & \emph{Application Programming Interface} (interfaz de programaci\'on de aplicaciones). Servicio HTTP REST en FastAPI que centraliza los c\'alculos del modelo. \\
CI & \emph{Continuous Integration} (integraci\'on continua). Flujo automatizado en GitHub Actions que ejecuta las pruebas, verifica la cobertura del 100\,\% y env\'ia el reporte a SonarQube Cloud. \\
D5 & Decisi\'on de dise\~no D5 del documento 04 del repositorio \parencite[doc.~04, \S~1]{axioma2026engine}: no modificar los par\'ametros por defecto del modelo sin datos experimentales para evitar ajustar los supuestos a las conjeturas. \\
HR & Humedad relativa del aire ambiente, expresada en porcentaje (\%). \\
ICBF & Instituto Colombiano de Bienestar Familiar. Entidad p\'ublica cuya nota de 2025 reporta la entrega mensual de $1000\;\text{L}$ por comunidad mediante carrotanque en La Guajira. \\
OMS & Organizaci\'on Mundial de la Salud. Define el est\'andar de acceso b\'asico de $20\;\text{L}$ por persona al d\'ia \parencite{howard2020domestic}. \\
R1--R6 & Reglas del motor de control en orden de prioridad (\cref{sec:motor_reglas}); su equivalencia con los identificadores del c\'odigo est\'a en el \cref{anexo:trazabilidad}. \\
TRL & \emph{Technology Readiness Level} (nivel de madurez tecnol\'ogica). El proyecto se ubica en TRL 2--3 (formulaci\'on del concepto y validaci\'on anal\'itica por simulaci\'on). \\
\bottomrule
\end{tabularx}
